% Responsible statements: C_F0068 Part A, C_F0079 v4, C_F0063, and C_F0024 with C_F0068 Part C.
\section{Introduction}\label{sec:introduction}

Evaluating a noncommutative rational function at a tuple $X$ of bounded operators can yield an unbounded affiliated operator.
The free-field realization theorem of Mai, Speicher and Yin \cite[Theorem 1.1]{MSY23} shows that evaluation at $X$ embeds the free field into the $*$-algebra of affiliated operators when $\Delta(X)=d$.
Under this condition, we first establish algebraic characterizations of rationality in $X$.
We then derive two analytic criteria: one using liberation dual systems and the other using dual systems.

Our question is motivated by Connes's rationality conjecture for free group algebras \cite[Section IV.5, Remark 3]{Con94}, a noncommutative analogue of Kronecker's criterion \cite{Kro81}.
In its operator formulation, this criterion characterizes rationality by a finite-rank commutator condition \cite[Section IV.3, Proposition 1]{Con94}.
For a free group, Connes observed that commutators with the operator arising from its Cayley tree have finite rank on the rational closure of the group algebra in the reduced group $C^*$-algebra,
and conjectured the converse.

Duchamp and Reutenauer \cite[Theorem 7]{DR97} proved this conjecture for a free group $\bF_d$ with $d$ generators ($d\geq1$),
and showed that the rational closure coincides with the division closure.
Their proof uses noncommutative rational series.
Using the bounded dual system for free Haar unitaries (see, for example, \cite[Proposition 6.3]{MSY23}),
one can express the Cayley-tree finite-rank condition as finite rank of the corresponding commutators.

Linnell \cite[Theorem 1.1]{Lin00} extends the free-group criterion to affiliated operators using left and right fraction representations.
Theorem~\ref{thm:analytic} extends this criterion to normal generators admitting a dual system in the sense of Definition~\ref{def:dual-system}, including free Haar unitaries with their bounded dual systems.
Linnell also identifies the bounded part of the affiliated division closure with the division closure in the reduced group $C^*$-algebra.
Together with this identity, Theorem~\ref{thm:analytic} recovers the criterion of Duchamp and Reutenauer.

Miyagawa \cite[Theorems 3.1 and 4.5]{Miy22} proves a semicircular analogue of these criteria, characterizing rationality by finite-rank commutators with the right annihilation operators (whose negatives form a bounded dual system in our convention).
His proof adapts the noncommutative rational-series approach of Duchamp and Reutenauer to the semicircular setting.
He also extends this characterization to affiliated operators using left and right fraction representations.

\subsection{A rationality criterion under maximal \texorpdfstring{$\Delta$}{Delta}}
Our first main result gives several equivalent characterizations under maximal $\Delta$ alone.
The quantity $\Delta$ was introduced by Connes and Shlyakhtenko in their paper \cite{CS05} on the homology of von Neumann algebras.

Throughout the paper, $(\M,\tau)$ is a tracial $W^*$-probability space, that is, a von Neumann algebra $\M$ equipped with a faithful normal tracial state $\tau$.
Let $X=(X_1,\ldots,X_d)\in\M^d$, with $d\geq1$.
Put $H:=L^2(\M,\tau)$ and $\Omega:=1$.
Inner products on $H$ are linear in the second variable, and we identify elements of $\M$ with left multiplication on $H$.
The condition $\Delta(X)=d$ means that there are no nonzero finite-rank operators $T_i\in\F(H)$ such that
\[\sum_i[\rho(X_i),T_i]=0;\]
see Section~\ref{subsec:delta-tensor-maps} for details.
Here $\rho(X_i)$ denotes the right multiplication operator by $X_i$ on $H$.

We write $L^0(\M)$ for the $*$-algebra of closed densely defined operators affiliated with $\M$.
We identify $H$ with the square-integrable affiliated operators in $L^0(\M)$ \cite[Theorem~7.3.2]{AP18}.
In particular, $H\otimes_{\C}H\subseteq L^0(\M)\otimes_{\C}L^0(\M)^{\mathrm{op}}$.
Write $\A_X:=\C\langle X_1,\ldots,X_d\rangle$ for the algebra generated by $X_1,\ldots,X_d$.
We denote its division closure in $L^0(\M)$ by $\Div_{L^0(\M)}\A_X$ and call an affiliated operator \emph{rational} when it belongs to this closure; see Definition~\ref{def:rationality}.

\begin{definition}\label{def:fraction-pair}
For $u\in L^0(\M)$, a \emph{fraction pair} is a quadruple $Q=(s,a,b,t)\in\M^4$ satisfying
\begin{equation}
 u=s^{-1}a=bt^{-1},\qquad s,t\text{ invertible in }L^0(\M).
 \label{eq:fractions}
\end{equation}
\end{definition}

\begin{definition}\label{def:fraction-commutator}
Let $D$ be a linear operator on $H$ with domain $\dom D$.
For a fraction pair $Q=(s,a,b,t)$ of an element $u\in L^0(\M)$, define the \emph{fraction commutator of $D$ and $Q$} by
\begin{equation}
 \Gamma[Q,D]:=sDb-aDt,
 \label{eq:fraction-commutator}
\end{equation}
with domain given by $\{\xi\in H: b\xi,t\xi\in\dom D\}$.
\end{definition}
The identities $ut=b$ and $su=a$ from \eqref{eq:fractions} give formally
\[
 s[D,u]t=sD(ut)-(su)D t=sD b-aD t=\Gamma[Q,D].
\]
That is, $s$ and $t$ remove the fractions in $u$ from the commutator expression.

Multiplication in $L^0(\M)\otimes_{\C}L^0(\M)^{\mathrm{op}}$ is given by
\[
 (u_1\otimes v_1^{\mathrm{op}})(u_2\otimes v_2^{\mathrm{op}})
 :=u_1u_2\otimes(v_2v_1)^{\mathrm{op}}.
\]
To state our first main result, we use the right multiplication maps in \cite[Section~3.1.1, equation~(3.1)]{CS05},
with the coefficient algebra enlarged from $\M\otimes_{\C}\M^{\mathrm{op}}$ to $L^0(\M)\otimes_{\C}L^0(\M)^{\mathrm{op}}$.

\begin{definition}\label{def:psi}
For $w\in L^0(\M)$, define the right multiplication map
\[
 \Psi_w:L^0(\M)\otimes_{\C}L^0(\M)^{\mathrm{op}}
 \longrightarrow L^0(\M)\otimes_{\C}L^0(\M)^{\mathrm{op}}
\]
by
\begin{equation}
 \Psi_w(k):=k(w\otimes1-1\otimes w^{\mathrm{op}})
 \label{eq:psi}
\end{equation}
\end{definition}

For $u,v\in L^0(\M)$, this gives
\[
 \Psi_w(u\otimes v^{\mathrm{op}})
 =uw\otimes v^{\mathrm{op}}-u\otimes(wv)^{\mathrm{op}}.
\]
For a tuple $X=(X_1,\ldots,X_d)\in\M^d$, we further write $\Psi_i:=\Psi_{X_i}$.

\Needspace{4\baselineskip}
\begin{theorem}\label{thm:delta}
Suppose $\Delta(X)=d$.
For $u\in L^0(\M)$ the following are equivalent:
\begin{enumerate}
\item $u\in\Div_{L^0(\M)}\A_X$.
\item\label{item:delta-tensors} There are $k_i\in L^0(\M)\otimes_{\C}L^0(\M)^{\mathrm{op}}$ such that
\begin{equation}
 \sum_{i=1}^d\Psi_i(k_i)=u\otimes1-1\otimes u^{\mathrm{op}}.
 \label{eq:tensor-criterion}
\end{equation}
\item\label{item:delta-all-fraction-tensors}
For every fraction pair $Q=(s,a,b,t)$ for $u$, there exist $\ell_1,\ldots,\ell_d\in L^0(\M)\otimes_{\C}L^0(\M)^{\mathrm{op}}$ such that
\begin{equation}\label{eq:all-fraction-tensor-criterion}
 \sum_{i=1}^d\Psi_i(\ell_i)
 =a\otimes t^{\mathrm{op}}-s\otimes b^{\mathrm{op}}.
\end{equation}
\item\label{item:delta-bounded-fraction-tensors}
There exist a fraction pair $Q=(s,a,b,t)$ for $u$ and tensors $\ell_1,\ldots,\ell_d\in\M\otimes_{\C}\M^{\mathrm{op}}$ satisfying \eqref{eq:all-fraction-tensor-criterion}.
\item\label{item:delta-finite-rank} There exist a fraction pair $Q=(s,a,b,t)$ for $u$ and a tuple $T=(T_1,\ldots,T_d)\in\F(H)^d$ such that
\begin{equation}
 \sum_{i=1}^d[\rho(X_i),T_i]=\Gamma[Q,P_\Omega].
 \label{eq:certificate}
\end{equation}
\end{enumerate}
In~\textup{(\ref{item:delta-tensors})}, $k_1,\ldots,k_d$ are uniquely determined by $u$.
The tensors $\ell_1,\ldots,\ell_d$ in \textup{(\ref{item:delta-all-fraction-tensors})} and \textup{(\ref{item:delta-bounded-fraction-tensors})} and the tuple $T$ in~\textup{(\ref{item:delta-finite-rank})} are uniquely determined by the fraction pair $Q$.
\end{theorem}

Here $P_\Omega$ denotes the orthogonal projection of $H$ onto $\C\Omega$.
By Definition~\ref{def:fraction-commutator},
\[
\Gamma[Q,P_\Omega]=sP_\Omega b-aP_\Omega t\in\F(H).
\]
The left-hand sides of equations~\eqref{eq:tensor-criterion} and~\eqref{eq:certificate} are values of degree-one Hochschild differentials with the corresponding coefficient bimodules; see Remark~\ref{rem:boundary-interpretation}.

\begin{remark}
When the equivalent assertions of Theorem~\ref{thm:delta} hold, the tensors in assertion~\textup{(\ref{item:delta-tensors})} are the free difference quotients $k_i=\partial_i u$ of $u$ with respect to $X_1,\ldots,X_d$;
see the implication $(1)\implies(2)$ in its proof.
The implication $(2)\implies(3)$ follows as $\Psi_i$ is a right multiplication map.
The implication $(3)\implies(4)$ is obtained by clearing denominators in $\ell_i$.
The implication $(4)\implies(5)$ is due to the identification of the tensors with finite-rank operators.
The key input for $(5)\implies(1)$ is that $\Div_{L^0(\M)}\A_X$ is a division ring, as guaranteed under $\Delta(X)=d$ by \cite[Theorem~1.1]{MSY23}.
\end{remark}

\subsection{Analytic criteria}
Condition~\textup{(\ref{item:delta-finite-rank})} in Theorem~\ref{thm:delta} requires a fraction pair $Q$ for which the commutator equation \eqref{eq:certificate} has the (unique) finite-rank solution $T\in\F(H)^d$.
The next theorem gives a rationality criterion that is valid for every fraction pair.
That is, \eqref{eq:certificate} admits a finite-rank solution either for every fraction pair or for none.

To formulate this analytic criterion, we write $\M\Omega:=\{x\Omega:x\in\M\}\subset H$.
We use the following variant of the (liberation) dual systems in \cite[Definition~13.1]{Voi99}.
\begin{definition}\label{def:liberation-dual-system}
Let $X=(X_1,\ldots,X_d)\in\M^d$.
A \emph{liberation dual system} for $X$ is a tuple $(L_1,\ldots,L_d)$ of linear operators on $H$ such that
\[
 \M\Omega\subseteq\dom L_i\cap\dom L_i^*
 \qquad(1\leq i\leq d),
\]
and
\[
 [X_j,L_i]\xi=\delta_{ij}[X_i,P_\Omega]\xi
 \qquad(\xi\in\M\Omega,\ 1\leq i,j\leq d).
\]
We call this system \emph{normalized} if
\begin{equation}\label{eq:liberation-normalization}
 \sum_{i=1}^dL_i\xi=P_\Omega\xi
 \qquad(\xi\in\M\Omega).
\end{equation}
\end{definition}
Every liberation dual system can be normalized, as shown in Lemma~\ref{lem:liberation-normalization}.

\Needspace{6\baselineskip}
\begin{theorem}\label{thm:intrinsic-analytic}
Let $X=(X_1,\ldots,X_d)\in\M^d$, and suppose that $X_i-\lambda$ is invertible in $L^0(\M)$ for every $1\leq i\leq d$ and every $\lambda\in\C$.
Let $(L_1,\ldots,L_d)$ be a normalized liberation dual system for $X$ in the sense of Definition~\ref{def:liberation-dual-system}.
Then, for every $u\in L^0(\M)$, the following are equivalent:
\begin{enumerate}
\item $u\in\Div_{L^0(\M)}\A_X$.
\item\label{item:intrinsic-analytic-existence-criterion}  There exist a fraction pair $Q$ for $u$ and $T_1,\ldots,T_d\in\F(H)$ such that
\begin{equation}
 \sum_{i=1}^d[\rho(X_i),T_i]=\Gamma[Q,P_\Omega].
 \label{eq:intrinsic-analytic-criterion}
\end{equation}
\item\label{item:intrinsic-analytic-criterion} For every fraction pair $Q$ for $u$, there exist $T_1,\ldots,T_d\in\F(H)$ satisfying \eqref{eq:intrinsic-analytic-criterion}.
\item\label{item:intrinsic-analytic-fraction-commutators} For every fraction pair $Q$ for $u$ and every $i$,
\[
 \overline{\Gamma[Q,L_i]}\in[\rho(X_i),\F(H)].
\]
\end{enumerate}
For each fixed fraction pair $Q$, the tuple $(T_1,\ldots,T_d)$ in assertions~\textup{(\ref{item:intrinsic-analytic-existence-criterion})} and~\textup{(\ref{item:intrinsic-analytic-criterion})} is unique.
\end{theorem}

Here $[\rho(X_i),\F(H)]:=\{[\rho(X_i),T]:T\in\F(H)\}$ and $\overline{\Gamma[Q,L_i]}$ denotes the closure of the closable operator $\Gamma[Q,L_i]$.
The new characterization in assertion~\textup{(\ref{item:intrinsic-analytic-fraction-commutators})} can be understood as follows.
For any fraction pair $Q$, assertion~\textup{(\ref{item:intrinsic-analytic-fraction-commutators})} provides $T_i\in\F(H)$ with $[\rho(X_i),T_i]=\overline{\Gamma[Q,L_i]}$.
Summing over $i$ and using the normalization in equation~\eqref{eq:liberation-normalization}, we obtain
\[
 \sum_i[\rho(X_i),T_i]
 =\sum_i\overline{\Gamma[Q,L_i]}
 =\Gamma[Q,P_\Omega],
\]
so $(T_1,\ldots,T_d)$ is the solution in assertion~\textup{(\ref{item:intrinsic-analytic-criterion})}.

Theorems~\ref{thm:delta} and~\ref{thm:intrinsic-analytic} do not require $X_1,\ldots,X_d$ to be normal.
For normal generators, we obtain two rationality criteria.
In Section~\ref{sec:conjugate-variable-criterion}, we apply Theorem~\ref{thm:intrinsic-analytic} to normal generators with nonatomic spectral distributions and $L^2$ liberation gradients.
The case of generators admitting conjugate variables follows as a corollary.

Our third main result further treats normal generators admitting a dual system and identifies the finite-rank solutions $T_i$ with the closures of the corresponding fraction commutators.
We use the following variant of the dual systems in \cite[Definition~5.1]{Voi98}.

\begin{definition}\label{def:dual-system}
Let $X=(X_1,\ldots,X_d)\in\M^d$.
A \emph{dual system} for $X$ is a tuple $(D_1,\ldots,D_d)$ of linear operators on $H$ such that
\begin{equation}\label{eq:analytic-hyp}
 \M\Omega\subseteq\dom D_i\cap\dom D_i^*,\qquad
 [X_j,D_i]|_{\M\Omega}=\delta_{ij}P_\Omega|_{\M\Omega}
 \qquad(1\leq i,j\leq d).
\end{equation}
\end{definition}

\Needspace{6\baselineskip}
\begin{theorem}\label{thm:analytic}
Suppose that $X_1,\ldots,X_d\in\M$ are normal operators generating $\M$, and let $(D_1,\ldots,D_d)$ be a dual system for $X$ in the sense of Definition~\ref{def:dual-system}.
For every $u\in L^0(\M)$ and every fraction pair $Q=(s,a,b,t)$ for $u$, the following are equivalent:
\begin{enumerate}
\item $u\in\Div_{L^0(\M)}\A_X$.
\item For every $i$, $\Gamma[Q,D_i]$ has finite rank.
\item For every $i$, $\overline{\Gamma[Q,D_i]}\in\F(H)$.
\end{enumerate}
When these equivalent conditions hold, the rank of each $\overline{\Gamma[Q,D_i]}$ is independent of $Q$.
\end{theorem}

We deduce Theorem~\ref{thm:analytic} from the conjugate-variable case of Corollary~\ref{prop:normal-conjugates}.
Lemma~\ref{lem:dual-conjugates} shows that normal generators admitting a dual system also admit conjugate variables.
The corollary then supplies finite-rank certificates for rationality, and Lemma~\ref{lem:finite-rank-identification} identifies their coordinates with the corresponding fraction-commutator closures.
For the converse, Lemma~\ref{lem:fundamental} constructs a normalized liberation dual system $(L_1,\ldots,L_d)$ from $(D_1,\ldots,D_d)$.
For this system,
\[
 \overline{\Gamma[Q,L_i]}
 =[\rho(X_i),\overline{\Gamma[Q,D_i]}]
\]
whenever $\overline{\Gamma[Q,D_i]}$ is bounded.
Consequently, when these closures have finite rank, they are the finite-rank solution required in Theorem~\ref{thm:intrinsic-analytic}\textup{(\ref{item:intrinsic-analytic-existence-criterion})}, i.e.,
\[
 \sum_{i=1}^d[\rho(X_i),\overline{\Gamma[Q,D_i]}]=\Gamma[Q,P_\Omega].
\]
The converse certificate criterion in Corollary~\ref{prop:normal-conjugates} now gives rationality.

A dual system for normal generators yields conjugate variables by Lemma~\ref{lem:dual-conjugates}; the arcsine tuple of Remark~\ref{rem:liberation-without-conjugates} has a normalized liberation dual system but no conjugate variable, hence no dual system.
The arcsine tuple therefore satisfies the hypotheses of Theorem~\ref{thm:intrinsic-analytic} but not those of Theorem~\ref{thm:analytic}.
Hence the hypotheses of Theorem~\ref{thm:analytic} are strictly stronger than those of Theorem~\ref{thm:intrinsic-analytic}.

At present, we do not know whether normal generators admitting conjugate variables necessarily admit a dual system in the sense of Definition~\ref{def:dual-system}.
We establish the conjugate-variable case before proving the dual-system criterion.
 
For normal operators admitting bounded dual systems, the fraction commutators are bounded and $\overline{\Gamma[Q,D_i]}=\Gamma[Q,D_i]$.
Theorem~\ref{thm:analytic} therefore recovers the criteria of Duchamp and Reutenauer and Linnell for free Haar unitaries, and of Miyagawa for free semicircular operators.

\subsection{Organization}\label{subsec:organization}

Section~\ref{sec:preliminaries} collects the algebraic preliminaries on division closures and the free field, followed by fractions of affiliated operators, finite-rank tensor realizations and the consequences of maximal $\Delta$.
Section~\ref{sec:maximal-delta} proves Theorem~\ref{thm:delta} using denominator clearing and the tensor realization of finite-rank operators.
Section~\ref{sec:liberation-dual-system-criterion} proves that liberation dual systems give maximal $\Delta$ when every $X_i-\lambda$ is regular, and establishes Theorem~\ref{thm:intrinsic-analytic} by recovering algebraic $L^2$ tensors.
Section~\ref{sec:conjugate-variable-criterion} constructs normalized liberation dual systems from $L^2$ liberation gradients and applies Theorem~\ref{thm:intrinsic-analytic} to normal generators with nonatomic spectral distributions.
It then obtains the conjugate-variable case as Corollary~\ref{prop:normal-conjugates} and gives an arcsine example separating the two hypotheses.
Section~\ref{sec:dual-system-criterion} proves independence of the fraction pair, identifies finite-rank certificates with fraction-commutator closures and establishes conjugate variables from a dual system.
It then deduces Theorem~\ref{thm:analytic} from Corollary~\ref{prop:normal-conjugates}, using the summed fraction-commutator identity for the converse.
Section~\ref{sec:examples} uses a circular operator to illustrate the role of normality in Theorem~\ref{thm:analytic}.
